\begin{solapartat}
\[ E_\textit{fotons} = W_0 + E_{c\ \textit{màxima}} \Rightarrow W_0 = \frac{hc}{\lambda} - E_{c\ \textit{màxima}} \quad\punts{0,2} \]
El potencial de frenada és una manera de determinar l'energia cinètica màxima:
\[ E_{c\ \textit{màxima}} = q\,V_\textit{potencial de frenada} \Rightarrow \]
\begin{align*}
W_0 &= \frac{hc}{\lambda} - q\,V_\textit{potencial de frenada} \quad\punts{0,2}\\
&= \frac{6,63\times 10^{-34}\cdot 3\times 10^{8}}{560\times 10^{-9}} - 1,6\times 10^{-19}\cdot 0,95 = 2,03\times 10^{-19}\ \mathrm{J} \quad\punts{0,2}
\end{align*}
On $W_0$ és la funció de treball.

La freqüència llindar per produir efecte fotoelèctric és:
\[ \nu_0 = \frac{W_0}{h} = \punts{0,2} = 3,06\times 10^{14}\ \mathrm{Hz} \quad\punts{0,2} \]
\end{solapartat}

\begin{solapartat}
\[ \text{Si } \lambda > \lambda_0 \Rightarrow \frac{1}{\lambda} < \frac{1}{\lambda_0} \Rightarrow E(\lambda) = \frac{hc}{\lambda} < \frac{hc}{\lambda_0} = W_0 \Rightarrow \]
No es pot produir efecte fotoelèctric \punts{0,5}
\[ \text{Si } \nu > \nu_0 \Rightarrow E(\nu) = h\nu > h\nu_0 = W_0 \Rightarrow \]
En aquest cas sí que es produirà efecte fotoelèctric \punts{0,5}
\end{solapartat}
